\section{Claim ledger and conclusion}
\label{sec:claim-boundary}

\subsection{Exact conclusions}

The paper proves:
\begin{enumerate}[label=\textup{(\arabic*)},leftmargin=2.8em]
 \item a literal divisor-poset zeta lift of every \(q\)-resolved primitive
       entry family;
 \item exact M\"obius inversion recovering every \(q\)-slice from the
       lifted diagonal profile;
 \item the diagonal filtration
       \((F^{\rm Z})^{[a]}=0\) for \(a\nmid|\Delta|\);
 \item an exact divisor-supported high-diagonal tail formula;
 \item layered actual-support bounds that retain cancellation;
 \item the complete affine class of provenance filters, containing both
       zero and zeta completions;
 \item an attained reduced primal and dual content-flow gauge;
 \item Hermitian reversal and independent-edge Perron factorization;
 \item an exact hierarchy from the signed norm through unrestricted,
       provenance, zero, and zeta majorants; and
 \item explicit examples proving that neither zero nor zeta completion
       universally dominates the other.
\end{enumerate}

\subsection{Statements not proved}

The paper does not prove:
\begin{itemize}
 \item that the zeta lift improves the zero completion on the actual
       growing carrier;
 \item that the provenance optimum has a mass-normalized saving;
 \item that divisor-supported diagonal sparsity makes the exact tail
       small;
 \item a complete row-sum or Perron decay estimate;
 \item a single filter uniform over a growing coefficient class;
 \item that an obstruction for provenance filters covers unrestricted
       analytic completions;
 \item that failure of an absolute majorant makes the signed matrix large;
 \item closure of any later physical reassembly gate;
 \item the strict \(1/400\) endpoint inequality;
 \item a result special to \(h_0=2\); or
 \item a parity-breaking, prime-pair, or twin-prime theorem.
\end{itemize}

\subsection{Audit levels}

\begin{center}
\begin{tabular}{p{0.30\textwidth}p{0.48\textwidth}l}
\toprule
Result & Scope & Level\\
\midrule
Zeta lift, inversion, diagonal filtration &
All finite divisor-resolved primitive kernels & L0\\
Provenance gauge and reduced dual &
All finite fixed-slice affine filter classes & L0\\
Hermitian/Perron factorization &
All finite independent unordered-edge classes & L0\\
Literal content specialization &
Complete TPC--74 fixed-\(h_0\) missing-high carrier & L1\\
Normalized provenance saving &
Complete growing physical row family & Open L2\\
Physical reassembly and endpoint &
One common carrier and coefficient scope & Open\\
\bottomrule
\end{tabular}
\end{center}

\subsection{Conclusion}

A \(q\)-faithful nonprimitive completion in this TPC chain is now exact.
It is the
divisor-poset zeta transform
\[
 F^{\rm Z}(du,dv)
 =
 \sum_{d\mid q\mid|\Delta|}L_q(u,v),
\]
and its full diagonal profile is invertible.  At the level of the
analytic profile, the lift therefore neither loses nor ambiguously
mixes the literal content data.  It does not create physical atoms:
it reorganizes the given amplitudes into nested common-factor layers,
all supported on divisors of the same fixed row-pair difference.

The zeta rule is not privileged analytically.  The broader
provenance-compatible gauge optimizes every legal content-flow weight and
has the exact dual condition
\[
 B_{\boldsymbol L}^*A_D^*z=0.
\]
This turns a vague search for a good \(q\)-compatible interpolation into
a matched fill/obstruction problem.  The comparison examples show why
both outcomes are needed: zeta can lower the complete cost, and it can
also be strictly worse than the zero completion.

The next bottleneck is no longer algebraic definition but sparse actual
incidence.  TPC--77 is planned to investigate decomposition of that
incidence into independently certifiable components.  Only after a real
carrier census identifies which components dominate will it be sensible
to spend a number-theory gate on their asymptotic mass.
